\documentclass[11pt,a4paper]{article}
\usepackage{amsmath,amssymb,amsthm}
\usepackage[margin=1.05in]{geometry}
\usepackage{longtable}
\usepackage{hyperref}

\theoremstyle{plain}
\newtheorem{theorem}{Theorem}
\newtheorem{conjecture}[theorem]{Conjecture}
\newtheorem{corollary}[theorem]{Corollary}
\theoremstyle{definition}
\newtheorem{remark}[theorem]{Remark}

\renewcommand{\SS}{\mathfrak{S}}

\title{A negative map of the Elliott--Halberstam route to Goldbach\\
\large with one unconditional theorem, one conjecture,\\
and eighteen pre-registered closures,\\
two of which this paper withdraws}

\author{[author]}
\date{Working paper, increment 200 of the \texttt{goldbach-ln2}
program\\
\small All measurements reproducible from the accompanying repository}

\begin{document}
\maketitle

\begin{abstract}
Huang and Li proved that binary Goldbach for large even $N$ follows
from $EH_\mu(N^{\theta'})$ for a single $\theta' > 1/2$. This paper
reports a four-day computational and analytic campaign that mapped
that hypothesis from both of its sides. The demand side closes at the
level of identities. The supply side does not close: we report the
negative map, and also the re-audit that withdrew two of its own
closures once the target was corrected.

On the \emph{demand} side we prove one unconditional theorem
(Theorem~\ref{thm:A}): the M\"obius-weighted, fixed-class correlation
sum with weight $w_k = 1$ is $\ll_A N(\log N)^{-A}$, so Huang--Li's
$E_4$ consumption of $EH_\mu$ is unnecessary and the whole demand
collapses to the single scalar $E_3$. We then show
(Theorem~\ref{thm:C}) that this scalar is \emph{equivalent} to binary
Goldbach, by an unconditional identity which is Huang--Li's own
equation (22); the root cause is $\mu * \log = \Lambda$. The
demand-side search for a weaker sufficient condition is therefore
closed at the level of identities. Net progress toward Goldbach: zero.

On the \emph{supply} side the campaign produced one formal conjecture
(Conjecture~\ref{conj:L}): every M\"obius family probed factorizes as
a deterministic local mask times an exactly Gaussian fluctuation, so
that every ``sub-random'' reading in the literature of this problem
that we could reproduce is mask accounting. What the Goldbach chain
consumes is the amplitude half of that conjecture, and we report
\textbf{eighteen pre-registered closures} --- five route adjudications
against the source papers' explicit lemma hypotheses, ten
kill-tested technique designs, and three representation-class
experiments --- each closed with numbers, together with the five
structural constraints they impose on any future technique. We then
\emph{withdraw two of those closures ourselves}: a correction to the
stated target invalidated every death sentence that had been passed on
a budget rather than a structural basis, and the re-audit of all
eighteen is reported in full.

The deliverable is the map, not a theorem toward Goldbach. We state
throughout what is and is not claimed, and we record thirty-one
corrections to our own earlier statements, including three refutations
of our own derivations and one false-positive trap that would have
produced an apparent refutation of $EH_\mu$.
\end{abstract}

\tableofcontents

\section{Introduction}

\subsection{The conditional frame}

Let $N$ be a large even integer,
$\tilde r(N) = \sum_{n<N}\Lambda(n)\Lambda(N-n)$, and $\SS(N)$ the
singular series. Write $EH_\mu(Q)$ for the assertion that for every
$A>0$
\[
  \sum_{q\le Q}\max_{y<N}\max_{(a,q)=1}
  \Bigl|\sum_{\substack{n\le y\\ n\equiv a\,(q)}}\Lambda(n)\mu(N-n)
   -\frac{1}{\varphi(q)}\sum_{n\le y}\Lambda(n)\mu(N-n)\Bigr|
  \ll_A \frac{N}{(\log N)^A}.
\]
Huang and Li~\cite{HL}, continuing Pan~\cite{Pan}, proved that
$EH(N^{\theta}(\log N)^{2A+8})$ together with $EH_\mu(N^{1-\theta})$
implies $\tilde r(N)\ge \SS(N)(1-A(N))N + O(N(\log N)^{-A})$, and
deduced (their Corollary~1) that in view of Bombieri--Vinogradov,
binary Goldbach for all sufficiently large even integers follows from
$EH_\mu(N^{\theta'})$ alone for any single $\theta' > 1/2$.

This is the sharpest conditional frame we are aware of: one sentence
stands between classical technology and Goldbach. The campaign
reported here asked what that sentence actually costs.

\subsection{The organizing principle: two couplings}

The hypothesis contains a product of two M\"obius factors, and there
are exactly two ways their arguments can be coupled. This turns out to
organize everything.

\begin{itemize}
\item \textbf{Divisibility-coupled} (the \emph{demand} side). Inside
  Huang--Li's proof, $EH_\mu$ is consumed only in the fixed residue
  class $n \equiv N \pmod k$, i.e.\ $k \mid N-n$: one M\"obius
  argument divides the other. Section~\ref{sec:demand}.
\item \textbf{Difference-coupled} (the \emph{supply} side). Any attempt
  to \emph{prove} $EH_\mu$ by dispersion arrives at the dilate field
  $D(k) = \sum_m \mu(m)\mu(N-mk)$, where the two arguments are related
  by a difference, not a divisibility. Section~\ref{sec:supply}.
\end{itemize}

\noindent
The two behave completely differently, and each is closed for its own
reason:

\begin{center}
\begin{tabular}{p{0.20\textwidth}p{0.34\textwidth}p{0.34\textwidth}}
\hline
 & \textbf{Divisibility-coupled} & \textbf{Difference-coupled}\\
\hline
Divisor switch & exact, and free & exact, but useless\\
M\"obius lands on & the \emph{short} variable, $m\le N^{1/2-\delta}$ &
 the \emph{long} variable, $m \asymp N/K \ge N^{2/3}$\\
Consequence & Bombieri--Vinogradov closes it: Theorem~\ref{thm:A} &
 no known machine applies\\
But & the weighted version is \emph{equivalent} to Goldbach
 (Theorem~\ref{thm:C}) & no named coupling surface found, internal or
 external --- with one route re-opened on re-audit
 (\S\ref{sec:closures}, \S\ref{sec:reaudit})\\
\hline
\end{tabular}
\end{center}

\subsection{What this paper claims}

\begin{itemize}
\item \textbf{Claimed}: Theorem~\ref{thm:A}, Corollary~\ref{cor:B} and
  Theorem~\ref{thm:C}, with full proofs in the companion note
  \cite{ThmA}; the defect in \cite{HL}'s equation (18) and its repair;
  all measurements, which are reproducible; the closure table of
  \S\ref{sec:closures}.
\item \textbf{Conjectured}: Conjecture~\ref{conj:L}.
\item \textbf{Not claimed}: any theorem toward Goldbach.
  Theorem~\ref{thm:A} is unconditional but Goldbach-neutral: it
  removes exactly the half of the demand that carries no Goldbach
  content.
\end{itemize}

\section{The demand side}\label{sec:demand}

\subsection{Where $EH_\mu$ is actually consumed}

For $(k,N)=1$ and $1\le t<N$ put
\[
  E_\mu(t;k) = \sum_{\substack{n\le t\\ n\equiv N\,(k)}}
    \Lambda(n)\mu(N-n) - \frac{1}{\varphi(k)}\sum_{n\le t}
    \Lambda(n)\mu(N-n).
\]
With $K = (N-1)/\alpha$, Huang--Li's two consumption sites are
\[
  E_3(\alpha)=\!\!\sum_{\substack{k<K\\(k,N)=1}}\!\!\mu(k)\log k\,
  E_\mu(N;k),
  \qquad
  E_4(\alpha)=\!\!\sum_{\substack{k<K\\(k,N)=1}}\!\!\mu(k)\,
  \bigl[\text{$E_\mu$ with the extra weight }\log(N-n)\bigr].
\]
They differ in exactly one respect: the weight $w_k$ is $\log k$ in
$E_3$ and $1$ in $E_4$ (the $\log(N-n)$ inside $E_4$ is
$k$-independent and comes off by partial summation). In both cases
Huang--Li discard the signs $\mu(k)$ on the first line by the triangle
inequality, and only then appeal to $EH_\mu$. Keeping the signs is
what the following results exploit.

\subsection{The theorem}

\begin{theorem}[$w_k = 1$]\label{thm:A}
Fix $\theta' \in (1/2,1)$ and put $K = N^{\theta'}$. Then for every
$A>0$,
\[
  \sup_{1\le t<N}\Bigl|
   \sum_{\substack{k<K\\(k,N)=1}}\mu(k)\,E_\mu(t;k)\Bigr|
  \;\ll_{A,\theta'}\; \frac{N}{(\log N)^A}.
\]
Every ingredient except Bombieri--Vinogradov contributes only
$O(Ne^{-c\sqrt{\log N}})$.
\end{theorem}

The mechanism is one line. The class $n \equiv N \pmod k$ is
$k \mid N-n$, so the $k$-sum is a divisor sum over $u = N-n$;
switching and writing $u = mk$ gives, for squarefree $u$,
$\mu(u)\mu(k) = \mu(m)\mu^2(k)$ --- the M\"obius on the long variable
$k$ squares itself away, and the surviving M\"obius sits on the short
variable $m < N^{1-\theta'} \le N^{1/2-\delta}$. That is the classical
``M\"obius on the short variable plus Bombieri--Vinogradov''
configuration. Proofs, and the four mandatory corrections that an
adversarial review imposed, are in \cite{ThmA}.

\begin{corollary}\label{cor:B}
$E_4(\alpha) \ll_A N(\log N)^{-A}$ unconditionally, via the exact
identity $E_4(\alpha) = \int_1^{N-1} T_1(t)\,dt/(N-t)$. Hence
Huang--Li's Lemma 4 is not needed and the entire $EH_\mu$ demand
collapses to the single scalar $E_3(\alpha)$.
\end{corollary}

\begin{theorem}[$w_k = \log k$: the permanent closure]\label{thm:C}
Unconditionally,
\[
  E_3(\alpha) = \sum_{n<N}\Lambda(n)\Lambda(N-n)
   - \SS(N)\Bigl(N - \sum_{n<N}\Lambda(n)\mu(N-n)\Bigr)
   + O_A\!\left(N(\log N)^{-A}\right),
\]
which is Huang--Li's own equation (22). Hence
$E_3(\alpha)\ll_A N(\log N)^{-A}$ --- the weakest form of the
hypothesis their argument needs --- is \emph{equivalent} to
$\tilde r(N)\sim\SS(N)N$.
\end{theorem}

\begin{remark}
The two theorems are the same computation with two weights. For
$w_k = \log k$ the complete divisor sum is $-\Lambda$ rather than an
indicator, because $\mu * \log = \Lambda$ --- and that identity is
where Huang--Li \emph{start} (their (10)). The $\log k$ weight is the
carrier of the Goldbach content, so every divisor switch must hand it
back. This is closure at the level of identities, not of estimates: no
choice of $\theta'$, truncation, or smoothing can evade it.
\end{remark}

\subsection{The interior of the weight space is empty}

Theorems~\ref{thm:A} and~\ref{thm:C} are the two endpoints of a design
space: the weight $w_k$ attached to $\mu(k)$ in the demand functional.
It is natural to look for a weight between them --- complete divisor
sum still cheap, main-term coefficient still of size $1$ --- that would
extract $C(N)=\sum_{n<N}\Lambda(n)\mu(N-n)$ from
Bombieri--Vinogradov alone. There is none.

Write $b := \mu*w$, so $w_k=\sum_{d\mid k}b_d$; every weight has this
form. For squarefree $u$ one has $\sum_{k\mid u}\mu(k)w_k=\mu(u)b_u$,
so the switch identity reads
$B_w\,C(N)=\sum_{u<N}\Lambda(N-u)\mu^2(u)b_u-\mathcal R_w$ with
$B_w=\sum_{k<K,(k,N)=1}\mu(k)w_k/\varphi(k)$. Two exact facts then
oppose each other:

\begin{itemize}
\item \emph{Extraction}:
  $B_w=\sum_{d<K}b_d\,\frac{\mu(d)}{\varphi(d)}\rho_{dN}(K/d)$ with
  $\rho \ll e^{-c\sqrt{\log x}}$, so $B_w$ is controlled by the part of
  $b$ sitting \emph{at} the truncation point $K=N^{\theta'}$.
\item \emph{BV-accessibility}: expanding $w$ inside the residual gives
  moduli $md$ with $m<N^{1-\theta'}$, so $b$ may only be supported on
  $d\le N^{\theta'-1/2-\delta}$.
\end{itemize}

\begin{theorem}[no weight extracts $C(N)$]\label{thm:D}
If $b=\mu*w$ is supported in $[1,N^{\theta'-1/2-\delta}]$ then
$\|b\|_1/|B_w| \gg \exp(c\sqrt{(1/2+\delta)\log N})$, and the switch
identity yields at best
$|C(N)|\ll \exp(c\sqrt{\tfrac12\log N})\,N(\log N)^{-A}$ --- no saving
of any power of $\log N$.
\end{theorem}

\begin{theorem}[the no-go survives $EH$]\label{thm:Dprime}
If $\Lambda$ has level of distribution $\theta_E<1$, the same argument
gives $\|b\|_1/|B_w| \gg \exp(c\sqrt{(1-\theta_E)\log N})$. Closing the
gap requires $\theta_E=1$: equidistribution of $\Lambda$ to moduli of
size $N$ itself, where each progression holds $O(1)$ terms.
\end{theorem}

\begin{proposition}[the circle method has zero margin]\label{prop:E}
Writing $C(N)=\int_0^1 S_\Lambda S_\mu e(-N\alpha)d\alpha$, both
standard estimates lie at or above the trivial bound $\psi(N)\sim N$:
Cauchy--Schwarz gives $\sim(6/\pi^2)^{1/2}N(\log N)^{1/2}$, exceeding
it by a growing factor; and any bound
$\sup_\alpha|S_\mu|\cdot\|S_\Lambda\|_1$ is at least
$\|S_\mu\|_2\|S_\Lambda\|_1\gg N$ by Parseval. The measured margin
$N/(\sup|S_\mu|\,\|S_\Lambda\|_1)$ is $0.168,0.175,0.158,0.152$ at
$N=2^{14},\dots,2^{20}$ --- below $1$ and decaying.
\end{proposition}

Davenport's uniform bound $S_\mu\ll_A N(\log N)^{-A}$ is useless here:
it saves against $N$, while the pairing needs the scale $N^{1/2}$,
which Parseval forbids. This is the parity obstruction in
circle-method language.

Theorem~\ref{thm:D} assumes $b$ is supported low enough for BV, which
excludes the most natural family: $w_k=f(\log k)$ with $f$ polynomial.
There $b=\mu*\log^D=\Lambda_D$, the generalised von Mangoldt function,
so the complete part splits by the number of prime factors of $u$ into
a Goldbach-type piece ($r=1$) and Chen-type pieces ($r\ge2$).

\begin{proposition}[polynomial weights]\label{prop:Dpp}
$D=0$ gives $B_w\ll e^{-c\sqrt{\log K}}$. For a monomial $x^D$,
$D\ge1$, every term of the complete part is nonnegative
($\Lambda\ge0$, $\Lambda_D\ge0$), so it is $\asymp N(\log N)^{D-1}$
with a fixed sign. Cancelling across monomials requires tuning against
the asymptotics of the $r=1$ piece, which at $D=1$ is the binary
Goldbach sum --- circular. No polynomial weight makes the complete
part a classically known object.
\end{proposition}

Measured at $N=10^6,4\cdot10^6,1.6\cdot10^7$: the two pieces of
$\mathrm{CP}_2$ are the same order, with ratio $0.771,0.790,0.810$
drifting towards $1$, and $\mathrm{CP}_2/(N\log N) = 2.886, 2.949,
2.997$. (This refuted our own first prediction, that the top-$r$ piece
would dominate by a power of $\log N$; the closure rests on
nonnegativity instead.) The one analytically canonical tuning,
$f(x)=x^2-2\gamma x$, which makes $b$ mean-zero by killing the pole of
$\zeta W$ at $s=1$, moves the complete part by about five percent:
mean-zero is not smallness, because it removes the average of $b$ and
not its correlation with $\Lambda(N-\cdot)$.

The loss exponent $\sqrt{(1/2)\log N}$ is the $\sqrt N$ barrier itself:
it is the gap between where a weight must live to see $C(N)$ and where
Bombieri--Vinogradov permits it to live. This is a no-go for one
precisely specified method --- divisor switching with BV as the only
input --- over that method's entire weight space; it is not an
obstruction to other methods. Its purpose is to close the design space
that Theorems~\ref{thm:A} and~\ref{thm:C} opened.

\subsection{A defect in the published paper}

Equation (18) of \cite{HL} replaces the $n$-dependent constraint
$k < (N-n)/\alpha$, present in the definition of $S_2(\alpha)$, by the
$n$-free $k < (N-1)/\alpha$. The omitted term is
\[
  \Delta = \sum_{2\le m\le\alpha}\mu(m)\log m
    \sum_{\substack{k<(N-1)/\alpha\\ (k,m)=1}}\mu^2(k)\Lambda(N-mk),
\]
whose trivial bound $\ll N(\log N)^2$ exceeds the target. It closes by
the same machinery --- M\"obius again on the short variable --- under
hypotheses already assumed: unconditionally by Bombieri--Vinogradov in
the Corollary-1 regime, and under the assumed $EH$ in the Theorem-1
regime. Theorem 1 and Corollary 1 of \cite{HL} stand as stated.

\section{The supply side}\label{sec:supply}

\subsection{The consumable}

The chain's supply-side consumable, distilled, is
\[
  \textbf{E1 (weak form).}\quad
  D(k) = \!\!\sum_{\sqrt N < m \le N/k}\!\! \mu(m)\mu(N-mk),
  \qquad
  \sum_{k\sim K}|D(k)|^2 \ll (\log N)^{-2A-2}\sum_{k\sim K} M_k^2,
\]
for $K \le N^{1/3}$ dyadic, where $M_k$ is the length of the $m$-range.
A fixed log-power saving is the currency; $o(1)$ and $\log\log$ savings
are not consumable.

The normalisation is $M_k^2$, i.e.\ the \emph{trivial} bound, not
$M_k$, the square-root scale: the wall is
$T_{II} = \sum_{k\sim K}b_kD(k)$ with $|b_k|\ll\log N$, needed at
$N(\log N)^{-A}$, and Cauchy--Schwarz in $k$ gives
$\sum_k|D(k)|^2 \ll N^2/(K(\log N)^{2A+2})$ with
$\sum_{k\sim K}M_k^2 \asymp N^2/K$. We record this because an earlier
version of this program's own adjudication printed $M_k$ here, which
demands \emph{better than square-root cancellation} and is refuted by
our own measurements ($\sum|D|^2/\sum M_k = 0.305,0.310,0.319$ at
$N=10^8$ against a demand of $8.8\cdot10^{-6}$). The corrected target
is cleared by square-root cancellation with margin
$(N/K)(\log N)^{-2A-2}\to\infty$, though that margin is invisible at
accessible $N$. No route verdict changes: all five are blocked
structurally, at named-lemma level, not on margin.

\subsection{Conjecture L}

\begin{conjecture}[the factorization law]\label{conj:L}
Each M\"obius family probed in this program --- prime-indexed dilate
pairs $C_{k,k'} = \sum_{p\sim P}\mu(N-pk)\mu(N-pk')$, integer-indexed
dilates and their pairs, and M\"obius sums over thin progressions ---
factorizes as
\[
  \text{field} \;=\; \mathbf M \times \mathbf G,
\]
where $\mathbf M$ is the deterministic local mask, computed exactly by
finite modular enumeration from the $v_q$-data of $(N,k,k')$, and
$\mathbf G$ is, on the surviving support, a fluctuation that is
exactly Gaussian at half-normal scale: no mean field, no class
structure, kurtosis $3$, Wishart-consistent pair spectrum.
\end{conjecture}

The evidence is summarized in Table~\ref{tab:L}; full stamps and
one-shot reproduction are in the repository.

\begin{table}[t]
\centering\small
\begin{tabular}{p{0.30\textwidth}p{0.62\textwidth}}
\hline
Level & Stamp\\
\hline
Pair statistics & free class $0.97$--$1.02$, kurtosis $2.99$--$3.03$
 (two $N$)\\
Exact cells & every viable $(v_2,v_3)$ cell $0.99$--$1.06$\\
Mask, blind & $\mathrm{corr}(\text{pred},\text{obs}) = 1.0000$,
 amplitude error $1.5\%$; annihilation fractions predicted exactly at
 a fresh $N$\\
Pair matrix & $\lambda_{\max}$ at the Wishart null dead center
 ($z = -0.19$)\\
E1 itself at $10^9$ & band ratios $0.966 / 0.950 / 0.922$ over
 $K \sim 10^3/3\cdot10^3/10^4$, no growth\\
E1 definitive grid & 4 $N$ $\times$ 2 bands $\times$ 300 $k$ with
 bootstrap SE: mean $z = +0.02$, 7/8 cells within $1.25\sigma$, the
 suspect cell settled by near-census\\
\hline
\end{tabular}
\caption{Measured support for Conjecture~\ref{conj:L}.}
\label{tab:L}
\end{table}

Conjecture~\ref{conj:L} is \emph{stronger} than the chain needs: E1
requires only square-root cancellation on average. What is missing is
therefore not knowledge of the structure --- the mask is an algorithm
and $\mathbf G$ is featureless in every measurement --- but a proof
technique for the amplitude of a featureless object.

\section{The negative map: eighteen closures}\label{sec:closures}

Every entry below was pre-registered: the decision rule was fixed in
writing before the computation ran, and a design that failed its rule
was closed in one paragraph and never re-litigated. Route
adjudications were done in fresh context against the source papers'
verbatim lemma hypotheses.

\subsection{Adjudication of existing machinery (5)}

\begin{longtable}{p{0.24\textwidth}p{0.10\textwidth}p{0.56\textwidth}}
\hline
Route & Verdict & Blocking coordinate\\
\hline
\endhead
MRT / Lichtman 2020, shift $\to$ dilate & Blocked &
 the orthogonality factorization needs a \emph{linear} pair
 constraint ($h = m-n$); the dilate constraint
 $m'u - mu' = N(m'-m)$ is bilinear. The $h$-average is a translation;
 the $k$-average is a dilation, with no diagonalizing character
 family.\\
Tao 2016 entropy decrement, $k$-averaged & Blocked $\times 3$ &
 (i) no $k$-analog of approximate affine invariance; (ii) the sampling
 prime enters the phase multiplicatively, so the sparsification fails
 and the surviving input is the target itself; (iii) the saving is
 triple-log, below spec.\\
Lichtman 2020 rerun in dilate coordinates & Blocked &
 the congruence coupling is a genuine isomorphism but powers only the
 typical-set restriction; the decoupling blocks at the same bilinear
 coordinate, and the uniformity slot would need an $EH_\mu$-grade
 input (circular).\\
Dirichlet-polynomial fourth moment / Perron & Blocked &
 every log-saving mean-value pillar assumes coefficient
 multiplicativity in its own variable; $\mu(N-u)$ has none. Unfolding
 by Perron costs $T \gtrsim N^{1-o(1)}$; opening the fourth moment
 reproduces the binary correlation.\\
Partial slices & Partial &
 the type-I slice is already consumed; the \emph{$N$-averaged}
 theorem is provable but lands in exceptional-set territory, which
 Huang--Li cannot consume; no nontrivial fixed-$N$ slice exists.\\
\hline
\end{longtable}

\noindent
\textbf{The common obstruction.} $\mu(m)\mu(N-mk)$ couples its
variables simultaneously through the product $mk$ and the difference
$N-mk$; the pair constraint is bilinear and is diagonalized by no
single character family, additive or multiplicative, while each
route's decisive lemma consumes exactly that diagonalization as a
hypothesis.

\subsection{Technique designs, kill-tested (9)}

\begin{longtable}{p{0.05\textwidth}p{0.30\textwidth}p{0.55\textwidth}}
\hline
\# & Design & Result (pre-registered rule applied)\\
\hline
\endhead
K1 & multiplicative Fej\'er kernel on the exact dilation ladder &
 \textbf{dead}: full 63-divisor orbit $R^2 = 0.466$. Half the field's
 energy is invisible to its entire multiplicative orbit.\\
K2 & manufactured pair congruence (determinant/Kloosterman) &
 \textbf{dead}: 0/10 flags over $h = 1\ldots210$; congruent pairs are
 statistically $h$-blind, so the manufactured congruence buys nothing
 against the splitting cost.\\
K3 & Wishart / operator moment method &
 \textbf{dead}: $\mathrm{tr}(M^2),\mathrm{tr}(M^3),\mathrm{tr}(M^4)$
 dead-center on the Wishart null ($z = -0.48/+0.42/-0.03$). No
 sub-Wishart surplus to fund the exchange.\\
K4 & $N$-average descent &
 \textbf{dead}: 0/8 flags; the dual field is $\delta$-blind and
 $m$-uncorrelated.\\
R1 & zero-spectrum visibility (explicit formula) &
 \textbf{dead}: $R^2_{\text{zeros}} = 0.2152$ vs random-frequency
 $0.2196\pm0.0055$. The $m^{i\gamma}$ coherence is completely
 scrambled by the pairing.\\
R2 & determinant / Kloosterman phase &
 \textbf{dead}: regression $R^2 = -0.0001/+0.0004$ against controls
 $\pm0.0002$ over 48{,}000 coprime pairs. The field is phase-blind to
 its own determinant structure.\\
R3 & character transform of the $k$-average &
 \textbf{dead by analysis}: for $K \le N^{1/3}$ the transform deposits
 every component into thin progressions (moduli $\ge N^{2/3}$);
 Parseval forbids a statistical gain.\\
R4 & divisor switch (the program's own newest mechanism) &
 \textbf{dead}: see \S\ref{sec:R4}.\\
R5 & the circle method applied directly to $C(N)$ &
 \textbf{dead, zero margin}: both bills lie at or above the trivial
 bound, and the cap on the pointwise route is Parseval
 ($\sup\ge\|\cdot\|_2$), which no technique can move
 (Proposition~\ref{prop:E}).\\
\hline
\end{longtable}

\subsection{Representation classes (3, plus one refutation)}

\begin{longtable}{p{0.10\textwidth}p{0.25\textwidth}p{0.55\textwidth}}
\hline
Class & Test & Result\\
\hline
\endhead
C-I abelian & rational-peak energy vs mask-null &
 \textbf{closed}, 0/6: the abelian spectrum is mask-exact.\\
C-II inverse domain & special-frequency excess in the
 modular-inverse re-indexing &
 \textbf{closed on verification}: fired at 4/300 with $z = +10.97$
 under 8-draw nulls, collapsed under 64-draw nulls
 ($+4.14/+2.41/+1.19/+0.30$), permutation null concurring, second-$N$
 replication 0/4.\\
C-IV manufactured modularity & approximate Fricke law for
 $\Phi(z) = \sum t_m e(mz)$ &
 \textbf{closed}, 0/6 levels. A true modular form drives the defect to
 $\approx 0$, so absence is meaningful.\\
C-III Motohashi type & (no finite test) &
 \textbf{open}. A full theory draft was written and is refuted, by two
 structural coordinates: a classification that hides a type-II region,
 and a dual object that is the output of no legitimate transform. Its
 other two coordinates were budget calls against a target since
 corrected and do not hold (\S\ref{sec:reaudit}). The draft is dead;
 the route is not.\\
\hline
\end{longtable}

\subsection{R4 in detail: the divisor switch does not
localize}\label{sec:R4}

Theorem~\ref{thm:A} gave the program a mechanism it did not have when
the earlier rounds were designed, so it was tested in its turn.
Applied to the dilate field over its \emph{full} ranges the switch
gives an exact identity
\[
  \sum_{k\ge1}\ \sum_{m:\,mk\le N-1}\mu(m)\mu(N-mk)
  = \sum_{u<N}\mu(N-u)\sum_{m\mid u}\mu(m)
  = \mu(N-1),
\]
verified by brute force at $N = 5000$ and $N = 20000$. This is perfect
cancellation: $O(1)$ for a double sum of $\sim N\log N$ terms, far
beyond square root, and it is the strongest cancellation this program
found anywhere.

E1 imposes two restrictions that make the inner divisor sum
incomplete: the type-II cut $m > \sqrt N$ and the dyadic band
$k \sim K$. The kill-test asked whether any of the cancellation
survives, by measuring block sums $S_B(j) = \sum_{k\in\text{block}}D(k)$
against the $B$-independence that Conjecture~\ref{conj:L} predicts.

\emph{Result: none of it survives.} On $8000$ values of $k$ --- the
entire band on which the type-II field is non-empty, since
$m > \sqrt N$ forces $k < \sqrt N$ --- the block ratios are flat
($B=8$: $0.958$ and $1.023$ at two $N$, against $B=1$ baselines
$0.980$ and $0.979$), and the sharper diagnostic, the lag-1
autocorrelation of $D(k)/\sqrt{\mathrm{supp}(k)}$, reads
$+0.0104$ and $+0.0127$ against a standard error of $0.0112$: dead
zero, and if anything mildly positive rather than the negative
coherence a surviving residue would require. An initial run on
$2048\ k$ had shown a same-direction deficit at both $N$; it was
settled as block-count standard error, in accordance with this
program's practice after an earlier $+3.7\sigma$ reading regressed to
noise under power.

\emph{The mirror.} Switching the banded $L^1$ sum gives
$\sum_{k\sim K}D_{\text{full}}(k)
 = \sum_{u<N}\mu(N-u)\sum_{m\mid u,\ u/2K<m\le u/K}\mu(m)$, where
$m \asymp N/K \ge N^{2/3}$: the surviving M\"obius sits on the
\emph{long} variable, the exact opposite of the assignment that makes
Theorem~\ref{thm:A} work. The switch is not a technique that happens
to fail on the supply side; its single requirement --- M\"obius on the
short variable --- is precisely what the type-II cut forbids.

\subsection{Re-audit: what the corrected target costs the map}
\label{sec:reaudit}

Correction to the target (\S\ref{sec:supply}) changes the yardstick,
so every death sentence passed on a magnitude or scale-mismatch basis
had to be re-examined. Under the corrected target the E1 route carries
a \textbf{slack of $N/K \ge N^{2/3}$}: a technique may be lossy by up
to that factor and still deliver. Under the printed target the
required bound sat \emph{below} what nature supplies, so the slack was
zero and any loss was fatal. That is what re-opens things.

\begin{itemize}
\item \textbf{One closure is void.} C-III's first kill coordinate
  reads, verbatim, ``E1 is an L$^2(k)$ statement at Gaussian
  (square-root) scale: it demands $|D(k)|^2 \approx M_k$'', and
  concludes that the tree's output
  $\sum|D|^2 \ll (N^2/K)(\log)^{-2C}$ is ``off target by
  $N/K \ge x^{2/3}$''. But the corrected target \emph{is}
  $\sum|D|^2 \ll (N^2/K)(\log N)^{-2A-2}$ --- the same form, met
  whenever $C \ge A+1$. The distance measured was the distance to a
  scale the chain does not require.
\item \textbf{A second is void}: C-III's budget coordinate opens with
  the words ``\emph{no-margin} arithmetic'', and the no-margin premise
  is the misstated target; a loss of $x^{1/3}$ sits comfortably inside
  a slack of $N^{2/3}$. Its structural half softens too --- ``absolute
  value summation reaches at best trivial scale'' is no longer fatal
  once the target sits exactly one log-power below trivial --- leaving
  only its last clause, shift-averaged binary $\mu\mu$ correlations at
  log-power strength, which is a \emph{named} open problem.
\item \textbf{One is downgraded}: the L$^2$-versus-signed-L$^1$
  objection. The chain's consumable is the \emph{signed} sum
  $\sum_k b_kD(k)$; the L$^2$ statement enters only through
  Cauchy--Schwarz, so it is sufficient and strictly stronger than
  needed --- the implication runs the safe way. The measured price of
  that step is $\asymp\sqrt K$ ($58\times$, $45\times$ at
  $K=10^3,3\cdot10^3$), which the margin absorbs. The corollary is
  worth more than the verdict: Cauchy--Schwarz discards exactly the
  sign structure of $b_k$, whose power R4 already exhibited (for
  $b_k\equiv1$ the unrestricted sum is exactly $\mu(N-1)$). This
  program has been aiming at something strictly harder than its own
  target.
\item \textbf{Two are moot}: the ``excess $K^{2/3}$'' and the
  conductor-inflation half. Their budgets are probably void by the
  same arithmetic, but they belong to a proof program that two
  structural items kill independently of any target.
\item \textbf{The rest stand}, being structural: violated lemma
  premises, absent pair congruences, illegitimate transforms, and
  kill-tests that measured \emph{no} signal rather than
  \emph{insufficient} signal. Proposition~\ref{prop:E} is also
  unaffected: on $C(N)$ itself the slack is only a log power.
\end{itemize}

The net effect is that C-III moves from ``dead even with a free
Lemma S'' back to ``an open route needing a correct construction''.
What it needs is now stated exactly: a legitimate transform in place
of the draft's hybrid object; a classification covering the type-II
region the draft's bookkeeping hid; and quantitative averaged Chowla
--- shift-averaged binary $\mu\mu$ correlations with a \emph{fixed}
log-power saving, against a best known of $(\log)^{1-c}$.

That third item changes the character of the obstruction, and the
change is the most consequential thing in this section. The
adjudication's central finding was that the \textbf{dilate} average
admits no diagonalizing character family; but the \textbf{shift}
average is precisely the home ground of the Matomäki--Radziwi\l\l--Tao
machinery. If a legitimate transform converting one into the other
exists, what remains is quantitative strength inside an active area
rather than the absence of any coupling surface. None of this is
progress toward Goldbach --- it is the withdrawal of over-claimed
refutations, and the route is open again needing exactly what nobody
can currently do.

We record the hazard by name: \emph{scale-normalisation drift},
writing a target at the square-root scale when the chain consumes it
at the trivial scale. This program has now made that mistake twice
(the SEAM formalisation, and the target itself), and both times it was
caught only by comparing a stated target against our own measurements
of the same quantity. That comparison is now mandatory before any
target is used to adjudicate anything.

\section{What the closures constrain}

Any transform or invariance $T$ that linearizes the pair constraint at
a cost of at most a log power must satisfy all five of the following.

\begin{enumerate}
\item \textbf{(K1)} $T$ cannot act through divisibility sub-sums: the
  multiplicative orbit spans only half the field's energy, so $T$ must
  couple to the $p$-coprime core directly.
\item \textbf{(K2)} $T$ cannot manufacture its congruence externally:
  collapse structure pays only when it arises inside an intrinsic
  average that is already present.
\item \textbf{(K3)} $T$ cannot bootstrap from the field's own moments:
  every moment consumes higher $\mu$-correlations and the field
  carries no sub-Wishart surplus.
\item \textbf{(K4)} $T$ cannot descend from the $N$-average: its
  linearization is load-bearing and the $k$-average holds no shadow
  of it.
\item \textbf{(R4)} $T$ cannot be the divisor switch or a relative:
  the switch's cancellation is a property of \emph{complete} divisor
  sums, and the type-II cut makes every divisor sum incomplete.
\end{enumerate}

\noindent
Together with the round-2 findings --- zeros invisible through the
pairing, characters merely relocating the difficulty into thin
progressions, determinant phases blind --- the field offers no
coupling surface in any direction that has an existing mathematical
name. In the automorphic world, shifted convolutions
$\sum a(n)a(n+h)$ are controlled because the coefficients $a(n)$
\emph{come with} a spectral representation; $\mu$ has none, and every
route above failed exactly where it tried to borrow one. The technique
to be created is a spectral (or equivalent structural) representation
of $\mu$-pairs themselves, not a projection onto existing spectra ---
and its construction is a purely theoretical act, beyond what
kill-testing can reach.

\section{Methodology, and the corrections}

Three rules were followed throughout, and they are the reason we
believe the negative results.

\begin{enumerate}
\item \textbf{Pre-registration.} Every design's decision rule was
  written before its computation ran. Two designs fired and were then
  retracted on verification (C-II at $z = +10.97$; a $z = 9$ spectral
  excess that turned out to be our own null-design category error ---
  an i.i.d.-entry Wigner null applied to a Gram matrix).
\item \textbf{Adversarial review in fresh context.} Three of this
  program's own derivations were refuted this way (the E1 gate
  arithmetic and role-exchange; the SEAM formalization, falsified by
  our own data; the C-III tree at four coordinates). Theorem~\ref{thm:A}
  is the first derivation to survive the same protocol, and it did so
  only after four mandatory corrections.
\item \textbf{Power before belief.} A $+3.7\sigma$ band elevation
  regressed to $1.078\pm0.072$ under bootstrap; a $-2.4\sigma$ dip
  resolved at four times the sample; the R4 deficit dissolved at four
  times the $k$-range.
\end{enumerate}

\noindent
\textbf{Twenty-eight corrections} are recorded in the repository. Two
deserve mention here because they are traps rather than slips.

\begin{itemize}
\item \emph{The $(q,N)>1$ main-term trap.} In the proof of
  Theorem~\ref{thm:A}, assigning a main term $T_m/\varphi(q)$ to a
  residue class with $(q,N)>1$ --- which is degenerate, its true
  contribution being $O(\log^2 N)$ --- shifts the density by exactly
  $N/\varphi(N)$. Carried into the $\log k$ branch, that error
  produces an apparent \emph{refutation of $EH_\mu$}. This is the most
  plausible false positive the campaign met; without adversarial
  review it would have been announced.
\item \emph{Eight-draw-null inflation.} Nulls estimated from too few
  draws inflate maxima across a family of 300 tests; the C-II ``hit''
  survived 8-draw nulls and died under 64-draw nulls plus replication.
\end{itemize}

\noindent
One correction is to this campaign's own increment-192 record: the
bound in Theorem~\ref{thm:A} was briefly claimed as
$Ne^{-c\sqrt{\log N}}$. That is wrong --- Bombieri--Vinogradov yields
only $N(\log N)^{-A}$ for each fixed $A$, its internal Siegel--Walfisz
range being $q \le (\log N)^B$ with $B$ fixed --- and the claim is
withdrawn here. The corollaries are unaffected.

\section{Reproducibility}

All measurements run on a laptop with Python and numpy. One-shot
verification of the core corpus: \texttt{python code/verify\_all.py}.
The Theorem~\ref{thm:A} checks are
\texttt{code/thmA\_*.py} (switching identity exact to $5\cdot10^{-10}$;
residual versus predicted main term agreeing to 1--4\%; the
load-bearing density identity verified in exact rational arithmetic
for all 243 squarefree $m<400$); the closure experiments are
\texttt{code/e1\_forge\_*.py} and \texttt{code/e1\_constr\_*.py}.

\section{Summary}

\begin{itemize}
\item The demand side of the Huang--Li reduction is closed: one half
  is now an unconditional theorem, the other half is equivalent to the
  conclusion.
\item The supply side is closed to measurement: the object is
  featureless in every direction with a name, and the eighteen
  pre-registered closures say precisely which directions were checked
  and why each fails.
\item Net progress toward Goldbach is zero. The contribution is the
  map, one unconditional theorem that removes a Goldbach-neutral
  demand, one conjecture that any future construction must reproduce,
  and one defect report for the source paper.
\end{itemize}

\begin{thebibliography}{9}
\bibitem{HL} Jing-Jing Huang and Huixi Li, \emph{On the connection
between the Goldbach conjecture and the Elliott--Halberstam
conjecture}, arXiv:2005.03811v2 [math.NT], 2022.
\bibitem{ThmA} \emph{An unconditional bound for a M\"obius-weighted
correlation sum in fixed residue classes}, companion note,
\texttt{paper/theorem\_A.tex} in this repository.
\bibitem{Pan} Cheng-Dong Pan, \emph{A new attempt on Goldbach
conjecture}, Chinese Ann. Math. \textbf{3} (1982), 555--560.
\bibitem{Tao} Terence Tao, \emph{The logarithmically averaged Chowla
and Elliott conjectures for two-point correlations}, Forum Math. Pi
\textbf{4} (2016).
\bibitem{Li20} Jared Duker Lichtman, \emph{Averages of the M\"obius
function on shifted primes}, arXiv:2009.08969.
\bibitem{Li23} Jared Duker Lichtman, \emph{Primes in arithmetic
progressions to large moduli, and shifted primes without large prime
factors}, arXiv:2309.08522.
\bibitem{MV17} M. Ram Murty and Akshaa Vatwani, \emph{Twin primes and
the parity problem}, J. Number Theory \textbf{180} (2017), 643--659.
\end{thebibliography}

\end{document}
